\nonstopmode 
\magnification=\magstep 1 
\input amstex
\redefine\rightheadline{\eightrm \folio \hfil  THE IDEAL GENERATION
CONJECTURE \hfil }
\documentstyle{amsppt}  
\loadbold \loadmsbm\loadeufm 
\topmatter 
\title The Ideal Generation Conjecture for  General Rational  projective
Curves \endtitle
\author Ferruccio Orecchia\endauthor
\address Ferruccio Orecchia, Dipartimento di Matematica e Applicazioni "R.
Caccioppoli", Complesso 
Universitario di Monte S. Angelo
- Via Cintia , 80126 Napoli-Italy \endaddress
\email ORECCHIA\@MATNA2.DMA.UNINA.IT.\endemail\
\thanks Work partially supported by 
MURST\endthanks 

\keywords Generic position, ordinary multiple subvarieties,
conductor\endkeywords
\subjclass 14Q05\endsubjclass
\abstract  We pose a conjecture for the expected number 
of generators of the ideal of the union $C$ of
$s$ general rational irreducible  curves
in $\Bbb P^r$.
By using the
computer we prove the conjecture for $C$ of low degree $d$ (e.g. if $s=1$
for
 $d\leq 80$ and if $s\leq 10$ for $d\leq 40$).
\endabstract
\endtopmatter

\document
\head 1. Introduction.\endhead 
Let $C=C_1\cup...\cup C_s\subset  \Bbb P^r_k$, $r\geq 3$ be a
non-degenerate 
curve  union 
of $s$, $s\geq 1$ rational irreducible smooth disjoint
curves $C_i$ of degree  $d_i$, $i=1,...,s$, over an algebraically closed
field $k$. Then the degree of 
$C$ is $d=\sum_{i=1}^s d_i$.\medskip
 The curve
$C$  has maximal rank if, for every integer $n\geq 1$, the natural
restriction map 
$$\rho(n): H^0(O_{\Bbb P^r_k}(n)) \to H^0(O_C(n))$$
has maximal rank 
i.e. $\rho(n)$ is injective or surjective.\medskip
 Let $\alpha=Min\{n\in \Bbb N\mid  $$n+r\choose r$$ > dn+s\}$\medskip  

A  curve $C$ of maximal rank is minimally generated if 
the natural map
$$\sigma(\alpha): H^0(I_C(\alpha)) \otimes H^0(O_{\Bbb P^r_k}(1)) \to
H^0(I_C(\alpha+1))$$
is of maximal rank.
Fixed $r,s,d_1,...,d_s$, if there exists a curve $C$ of maximal rank
(respectively
of maximal rank and minimally generated), then, for 
general $C_1,...,C_s$, $C$ has maximal rank (respectively has maximal rank
and is minimally generated) [Corollary 2.6]. 

The problem of finding the values of $r,s,d_1,...,d_s$ for which there
exists
a curve $C$ of
 maximal rank has been completely solved. In fact, starting from work of
Hartshorne
and Hirshowitz
it has been shown in [4]  that if $r\geq 4$ then, for any $s,d_1,...,d_s$,
 there exists a curve $C$  maximal rank.
 Furthermore in [1] it has been shown that, if $r=3$,
there exists $C$ of maximal rank, for any $s,d_1,...,d_s$, except the
following 
five cases: $s=2, d_1=d_2=2;$   $s=2, d_1=2, d_2=4;$  $s=3, 
d_1=d_2=d_3=2;$  $s=3, d_1=d_2=2,d_3=4;$   $s=4, d_1=d_2=d_3=d_4=2$.

 The problem
of finding the values of $r,s,d_1,...,d_s$ for which there exists a curve
 $C$  of maximal rank and 
minimally generated is, in general, unsolved. In the case of lines, i.e.
$d_1=...=d_s=1$, of $\Bbb P^3$, it has been shown in [9] that, for any
$s\neq 4$,
there exist $s$ lines with maximal rank and minimally generated. In the
case $s=1$,
 recently in [3] and [10]
 it has been shown that for any
$d\geq 932$ or $d\leq 73$, ($d\neq 5$) 
there exists an irreducible rational smooth curve of $\Bbb P^3$ which has
maximal
rank and  is
minimally generated.




In this paper we pose the following conjecture, which we call  {\sl
 ideal generation conjecture}
(I.G.C. for short) for rational curves.

For any 
$r,s,d_1,...,d_s$, 
there exists a curve $C$ of maximal rank and minimally generated,
 except the following s-uples $(d_1,...,d_s)$.\medskip

If $r=3:$ 
$(5)$ (i.e. $s=1$ and $d_1=5$); $(1,2)$; $(2,2)$; $(2,3)$; $(2,4)$;
 $(3,4)$; $(1,1,2)$; 

$(1,2,2)$; 
$(2,2,2)$; $(2,2,3)$; $(2,2,4)$; $(1,1,1,1)$;
$(1,2,2,2)$; $(2,2,2,2)$; $(1,2,2,2,2)$.\medskip

If $r=4:$ 
$(5)$; $(2,3)$; $(1,1,2)$.
\medskip
If $r=5:$
$(3,4)$; $(2,2,2)$;\medskip
 
In particular, if $s=1$, the conjecture is that, for any $d$ and $r$
there exists an irreducible rational smooth curve of degree $d$ in $\Bbb
P^r$ 
with maximal rank
and minimally generated except when $d=5$ in $\Bbb P^3$ or $\Bbb P^4$.
This, 
if proved,
 would
 extend
to $\Bbb P^r$ the result of [3] and [10] in $\Bbb P^3$. 
Moreover, in the case of lines, if the 
conjecture is proved, we have that for any $s$ and $r$ there exist $s$
 lines in $\Bbb P^r$ with maximal rank and minimally generated except the
case of 
$4$ lines
in $\Bbb P^3$, thus extending 
the result of [9] in $\Bbb P^3$.

The conjecture has been formulated 
with the help of the computer.
In fact we show how to check the conjecture by computer for fixed
$r,s,d_1,...,d_s$
 using an algorithm [Algorithm 3.5] that reduces
the computation of the generators of the ideal of the curve $C$
to the computation of the generators
for the ideal of a set of points suitably chosen on $C$. This has allowed
us to prove the conjecture 
on our Personal Computer for $d\leq 80$, if $s=1$, and for $d\leq 40$, if
$s\leq 10$,
 but enlarging this bound is only a question of having
time and/or more powerful computers [Remark 3.10].

Note that the usual systems for computations in Commutative Algebra and 
Algebraic Geometry like CoCoA [5] seem less suitable for this
particular kind of computation [Remark 3.10].

This paper has been inspired by the ideal generation conjecture for a
finite set of points 
in general position, which was posed in [6] and 
although proved in many cases (see [2] and [8]), is still
open in its generality.

We would like to thank G. Albano and F. Cioffi for having done the
implementation 
of  Algorithm 3.5 of this paper in the language $C++$ and L. Chiantini for
many
useful discussions on this subject.
\head 1. Maximal rank.\endhead
In all the paper $k$ is a fixed algebraically closed field.\medskip
We recall that a coherent sheaf $\Cal F$ on $\Bbb P^r$ is  $n$-regular if
$H^i(\Cal F(n-i))=0$ for $i>0$. If $\Cal F$ is $n$-regular $\Cal F$ is
$n+1$-regular [11, Lecture 14].
\definition{Definition 1.1} Let $Y$ be a closed subvariety (also
reducible)
 of $\Bbb P^r$ and
$\rho(n): H^0(\Cal O_{\Bbb P^r_k}(n)) \to H^0(\Cal O_Y(n))$ be
the natural restriction map. We say that $Y$ has {\sl maximal rank} if,
for every
integer $n\geq 0$, 
 $\rho(n)$ has maximal rank
as a map of vector spaces i.e. it is injective or
surjective.\enddefinition

Let $I(Y)=\bigoplus_{n\geq 0}I(Y)_n\subset k[X_0,...,X_r]=R$ be the ideal
of 
$Y$ and $A=R/I(Y)$ be the coordinate
ring of $Y$. 

Set $H_Y(n)=dim_kA_n=dim_kR_n-dim_kI(Y)_n= $$n+r\choose r$$-dim_kI(Y)_n$
for the Hilbert function of  $Y$ and $P_Y(n)$ for the Hilbert polynomial
of $Y$.
 
In the following we say that {\sl $Y$ is generated in degree $n$} if 
the ideal $I(Y)$ of $Y$ can be generated by forms of degree $\leq n$.

\proclaim{Lemma 1.2} Let $Y$ be a curve and $n\geq 0$ a given integer.
 $Y$ is $n$-regular if and only if $\Cal O_Y(n-2)$ is non-special and
$\rho(n-1): H^0(\Cal O_{\Bbb P^r_k}(n-1)) \to H^0(\Cal O_Y(n-1))$ is
surjective in which case
$Y$ is generated in degree $n$.\endproclaim
\demo{Proof} [7, section 1]\qed\enddemo

 



\proclaim{Proposition 1.3} Let
 $C=C_1\cup...\cup C_s\subset  \Bbb P^r_k$, $r\geq 3$ be the union 
of  $s$ rational irreducible smooth disjoint 
curves $C_i$ of degree $d_i$, $i=1,...,s$.
 Denote with $d$  the degree $\sum_{i=1}^s d_i$ of $C$. Then:\medskip
 
(a) $H_C(n)\leq P_C(n)$ for any n;\medskip

(b) if $H_C(m)=P_C(m)$ for some integer $m$ then $C$ is $m+1$-regular,
whence
  the ideal $I(C)$ of $C$
is generated in degree $\leq m+1$ and
 $H_C(n)=P_C(n)$ for any $n\geq m$;
\medskip
(c) $P_C(n)=dn+s$; \medskip

(d) $H_C(n)\leq Min\{$$n+r\choose r$$, dn+s\}$ and
 $C$ has maximal rank if and only if
$H_C(n)=Min\{ $$n+r\choose r$$, dn+s \}$ \medskip
\endproclaim
\demo {Proof}
 $(a)$, $(b)$ Since $C$ is rational the line bundle $\Cal O_C(n)$ is
non-special
for any $n$. 
Then $C$ is $n$-regular if and only if $\rho(n-1)$ is surjective [Lemma
1.2].
Moreover, for any $n$,  $h^1(\Cal O_C(n))=0$ and then  $h^0(\Cal
O_C(n))=P(n)$.
Hence, for any $n$, $H(A,n)=dim_k(Im_{\rho(n)})\leq h^0(\Cal O_C(n))=P(n)$
 (which gives (a)) and $H_C(m)=P_C(m)$, for some integer $m$, if and only
if
 $C$ is $m+1$-regular.
 This gives $(b)$ in virtue of Lemma 1.2 and the fact that if $C$ is
$m$-regular,
for some integer $m$, then $C$ is $m+1$-regular.\medskip
 

$(c)$ Since the irreducible curves $C_i$, $i=1,...,s$ are smooth, rational
of degree $d_i$ it is well known that $P_{C_i}(n)=d_in+1$. Moreover the
$C_i$ are disjoint then
$P_C(n)=\sum_{i=1}^s P_{C_i}(n)=\sum_{i=1}^s (d_in+1)=dn+s$.
\medskip
$(d)$ The first claim follows from $(a)$ and $(c)$ since
 $H_C(n) \leq $$n+r\choose r$$ $
for any $n$.
The second claim follows by observing that $\rho(n)$ is injective if and
only if
$H_C(n)= $$n+r\choose r$$ $ and that $\rho(n)$ is surjective if 
and only if $H_C(n)=P_C(n)=dn+s$ by the previous
considerations.\qed\enddemo

In general the Hilbert function of a rational curve (not smooth) is not
bounded by 
the corresponding Hilbert polynomial as the following example shows.


\proclaim {Example 1.4} Let $x=u(u^5-v^5),y=v(u^5-v^5),z=v^6$
 be a parametrization of the rational sextic
$C:y^6-x^5z+y^5z=0$. Then it is easily seen that $P_C(n)=6n-9$
 while $H_C(3)=10> P_C(3)=9$.
\endproclaim

\proclaim{Proposition 1.5} Let $C=$$\bigcup_{i=1}^s$$ C_i  \subset \Bbb
P^r$
 be a non-degenerate curve,  
such that $C_i$, are rational irreducible curves parametrized by  maps 
$\Phi_i:\Bbb P^1 \rightarrow \Bbb P^r$
given by forms
of degree $d_i$, for any $i$. Let $d=\sum_{i=1}^s d_i$
and assume that $H_C(m)=dm+s$ for some $m$. Then  $C_i$ are smooth
disjoint curves of 
degree $d_i$ and $C$ is $m+1$-regular.\endproclaim
\demo{Proof} It is enough to prove that the curves $C_i$ are smooth and
disjoint
since then $C$ is $m+1$-regular by Proposition 1.3.\medskip
First we consider the case $i=1$.\medskip 

Let $C'$ be the $n$-uple
embedding  of $C$. clearly $C'\subset \Bbb P^{nd}$ and $H_C(m)=dm+1$ is
the
Hilbert function of the coordinate ring of $C'$ in  degree one. Then $C'$
is non 
degenerate. Hence $C'$ is the rational normal curve of degree
$nd$ in $\Bbb P^{nd}$. Since $C'$ is smooth also $C$ is smooth and has
degree $d$.\medskip

Let now $i>1$. \medskip
If $ C_i^*$ is the normalization of $C_i$ then $H_{C_i}(n)\leq
H_{C_i^*}(n) \leq d_in+1$, for any $n$ [Proposition 1.3]. Then 
$md+s=H_C(m)\leq \sum_1^s H_{C_i}(m)\leq \sum_1^s(d_im+1)=dm+s$ and then 
$H_{C_i}(m) = d_im+1$ for any $i=1,...,s$. Hence the curves 
$C_i$ are smooth of degree $d_i$ and $C$ has degree $d$. Now suppose that 
$C_i$ and $C_j$ are not disjoint for some $i,j$. It is easily seen that in
this case
$H_C(m)< \sum_1^sH_{C_i}(m)=\sum_1^s(d_im+1)=dm+s$ and this contradicts
the
hypothesis that $H_C(m)=dm+s$.
 \qed\enddemo
\proclaim{Theorem 1.6}  Let $C=$$\bigcup_{i=1}^s$$ C_i  \subset \Bbb P^r$
 be a  non-degenerate curve,  
such that $C_i$, are rational irreducible curves parametrized by  maps 
$\Phi_i:\Bbb P^1 \rightarrow \Bbb P^r$ 
given by forms
of degree $d_i$, for any $i$. Let $d=\sum_{i=1}^s d_i$ and
$\alpha=Min\{n\in\Bbb N\mid $$n+r\choose r$$ > dn+s\}$. Then the following
conditions are equivalent:\medskip
$($a$)$ the curves $C_i$ are smooth of degree $d_i$ and disjoint for any 
$i$ and $C$ has maximal rank; \medskip 
$($b$)$ $H_C(\alpha -1)=$$\alpha+r-1\choose r$$ $ and $H_C(\alpha)=d\alpha
+s$;\medskip
and imply that $I(C)$ can be generated by forms of degree $\alpha$ and
$\alpha+1$.\endproclaim
\demo{Proof} If $H_C(\alpha -1)=$$\alpha+r-1\choose r$$ $ then 
$H_C(n)=$$n+r\choose r$$ $ for any $n<\alpha$. Hence  the claim follows
from Propositions
$1.3$ and $1.5$\qed\enddemo
\head 2. Minimal generation.\endhead

Let $C=$$\bigcup_{i=1}^s$$ C_i  \subset \Bbb P^r$, $r\geq 3$,
be the union of $s$ rational irreducible smooth disjoint curves and denote
with
$\nu (C)$ the minimal number of generators of the ideal $I(C)$ of $C$. 
With $I_C$ we denote the
ideal sheaf of $C$.
Assume that $C$ has maximal rank and let
$\alpha=Min\{n\in\Bbb N\mid $$n+r\choose r$$ > dn+s\}$; 
By Theorem
1.6 the ideal 
$I(C)=\bigoplus_{n\geq 0}I(C)_n$ of $C$ can be generated by forms of
 degree  $\alpha$ and $\alpha+1$. Then if we set
 $$W_{\alpha+1}=X_0I(C)_\alpha+...+X_rI(C)_\alpha\subset I(C)_{\alpha+1}$$
we have $$\nu (I)=dim_kI(C)_{\alpha}+
dim_kI(C)_{\alpha+1}-dim_kW_{\alpha+1}$$
Furthermore:\medskip

(1) $dim_k W_{\alpha+1} \leq Min\{(r+1)dim_k I(C)_\alpha,dim_k
I(C)_{\alpha+1}\}$
and then:
$$ \nu (I(C)))  \geq 
dim_k I(C)_\alpha+ dim_k I(C)_{\alpha+1}-Min\{(r+1)dim_k I(C)_\alpha,dim_k
I(C)_{\alpha+1}\}$$
\definition { Definition 2.1}   $C$ of maximal rank
is minimally generated if:

$$ \nu (I(C))) =
dim_kI(C)_\alpha+
dim_kI(C)_{\alpha+1}-Min\{(r+1)dim_kI(C)_\alpha,dim_kI(C)_{\alpha+1}\}$$
that is
$$dim_kW_{\alpha+1}=Min\{(r+1)dim_kI(C)_\alpha,dim_kI(C)_{\alpha+1}\}$$
that is 
$$\sigma(\alpha): H^0(I_C(\alpha))
 \otimes H^0(O_{\Bbb P^r_k}(1)) \to H^0(I_C(\alpha+1))$$
is of maximal rank. \enddefinition

We want to show that maximal rank and minimal generation can be checked by
computing the generators of a suitably chosen set of points on
$C$.\medskip

\proclaim {Lemma 2.2} Let $C=$$\bigcup_{i=1}^s$$ C_i  \subset \Bbb P_k^r$
 be a  non-degenerate  curve
such that the $C_i$ are rational irreducible curves parametrized by maps 
$\Phi_i:\Bbb P_k^1 \rightarrow \Bbb P_k^r$,
 $\Phi_i([t_1,t_2])=[f_{i0}(t_1,t_2),...,f_{in}(t_1,t_2)]$, 
where $f_{ij}\in k[t_1,t_2]$
are homogeneous polynomials of the same degree $d_i$.
Let $m>1$ a fixed integer. Consider the set $V_i$ consisting of 
$h_i=r_im+1$ distinct
 points  of $C_i$, for any $i$, and 
let $V=\cup_{i=1}^s V_i$ be the set of all these $(h=\sum_{i=1}^s h_i)$
points and 
$I(V)= \oplus_{n\geq 0} I(V)_n$
be the ideal of $V$.
 Then $I(C)_n=I(V)_n$
for any $n\leq m$.\endproclaim
\demo{Proof} Clearly $I(C)_m\subset I(V)_m$.Then we have  to prove the
opposite 
inclusion.
Since $I(V)= \cap_{i=1}^s I(V_i)$  and  $I(C)= \cap_{n=1}^s I(C_i)$, 
it is enough to prove the claim in the case 
of one irreducible curve $C_i$. Let $P_1,...,P_{h_i}$, 
 be the points of $V_i$.
 If $f\in I(V_i))_n$ then
$f(\Phi_i(P_u))=f(f_{i0}(P_u),...,f_{ir}(P_u))=0$, for any $u=1,...,h_i$.
 Thus $f(f_{i0},...,f_{ir})$ is a
polynomial of degree $nd_i\leq md_i$ vanishing on the 
$md _i+1$  points $P_i$ and then
$f(f_1,...,f_r)=0$. But $C_i$ is the smallest variety containing
the points $[f_{i0}(a_0,a_1),...,f_{ir}(a_0,a_1)]$, for any 
$[a_0,a_1]\in \Bbb P^1$
and then $C_i \subset V(f)$
that is $f\in I(C_i)$.\qed\enddemo
 
 

\proclaim {Remark 2.3} The $h_i$ points of Lemma 2.2 are immediately found
giving  values $n.1_k$, $n\in \Bbb Z$, to the
two parameters $t_1,t_2$.\endproclaim

\proclaim{Theorem 2.4} Let $C=$$\bigcup_{i=1}^s$$ C_i  \subset \Bbb P_k^r$
 be a non-degenerate curve
such that $C_i$ are rational irreducible curves parametrized by maps 
$\Phi_i:\Bbb P_k^1 \rightarrow \Bbb P_k^r$
given by forms of degree $d_i$ , for any $i$. Let $d=\sum_{i=1}^s d_i$ and 
$\alpha=Min\{n\in\Bbb N\mid $$n+r\choose r$$ > dn+s\}$
Consider $h_i=d_i(\alpha+1)+1$ points of $C_i$ for any $i$.
Let $V$ be the set of all these $h=\sum_{i=1}^s h_i=d(\alpha+1)+1$ points
and let 
$I(V)= \oplus_{n\geq 0} I(V)_n$
 be the ideal of $V$. Set:

 $$W_{\alpha+1}=X_0I(V)_\alpha+...+X_rI(V)_\alpha\subset
I(V)_{\alpha+1}$$\medskip

Then $C$
has maximal rank  if and only if\medskip 

$H_V (\alpha -1)=$$\alpha+r-1\choose r$$ $  , $H_V(\alpha)=d\alpha
+s$\medskip

Moreover $C$ is  minimally generated if and only if
$$dim_kW_{\alpha+1}=Min\{(r+1)dim_kI(V)_\alpha,dim_kI(V)_{\alpha+1}\}$$\endproclaim 
\demo{Proof} Follows immediately from Theorem 1.6, Definition 2.1 and
Lemma 2.2.\qed\enddemo

let $V=\{P_1,...,P_h\}\in \Bbb P_k^r$  be a set of points .
 The vector space $I(V)_n$ is easily given by the null
space of a matrix with elements in $k$. In fact if 
$R_n=\{f\in k[X_0,...,X_r]\mid deg(f)=n \}$
then:\medskip


(2) $I(V)_n=\{f\in R_n\mid f(P_i)=0,$ $i=1,...,h \}$\medskip

If we denote with $T_i$, $i=1,...,u$, the terms of degree $n$ in the 
indeterminates $X_0,...,X_r$ then $S=\{T_1,...,T_u\}$ is a basis of the
$k$-vector space
$R_n$. Let $G_n(V)$ be the following $ $$n+r\choose r$$\times h$
matrix:\medskip

(3) $G_n(V)=(b_{ij})$ where $b_{ij}=T_i(P_j)$\medskip

If $f=a_1T_1+...+a_uT_u\in k[X_0,...,X_r]$, we set


 $$(f)_S=\pmatrix a_1\\a_2\\ \vdots\\a_h\endpmatrix$$


Then:\medskip

(4) $I(V)_n=\{f\in R_n\mid G_n(V)(f)_S=0\}$ \medskip

   that is $f\in I(V)_n$ if and only if $(f)_S$
 is a
vector of the null space $ N_n(V)$ of the matrix $G_n(V)$. This gives, by 
elementary linear algebra that:
\medskip
(5) $dim_k(I(V)_n= $$n+r\choose r$$- rk(G_n(V))$ then 
 $rk(G_n(V))=H_V(n)$ $(rk=rank)$.
\medskip
Let
 $\{(g_1)_S,...,(g_e)_S\}$
be a basis of $N_n(V)$, $e= $$n+r\choose r$$ -rk(G_n(V))$. then the
$e(r+1)$ polynomials $g_iX_j$ of $I(V)_{\alpha+1}$
span the vector space $W_{\alpha+1}$. Let:\medskip

(6) $M_{\alpha+1}(V)$ be the matrix
whose columns are $(g_iX_j)_S$, for any $i,j$\medskip
 then\medskip

(7) $dim_k W_{\alpha+1}=rk(M_{\alpha+1}(V))$ \medskip


\proclaim {Corollary 2.5} Let $C=$$\bigcup_{i=1}^s$$ C_i  \subset \Bbb
P_k^r$
 be a non-degenerate curve
such that $C_i$ are rational irreducible curves parametrized by maps 
$\Phi_i:\Bbb P_k^1 \rightarrow \Bbb P_k^r$
given by forms of degree $d_i$ , for any $i$. Let $d=\sum_{i=1}^s d_i$ and 
$\alpha=Min\{n\in\Bbb N\mid $$n+r\choose r$$ > dn+s\}$
Consider $h_i=d_i(\alpha+1)+1$ points of $C_i$ for any $i$.
Let $V$ be the set of all these $h=\sum_{i=1}^s h_i=d(\alpha+1)+1$ points.
Then
  the curves $C_i$ are smooth and disjoint and $C$ has maximal rank 
if and only if\medskip

$ rk(G_{\alpha-1}(V))=$$\alpha+r-1\choose r$$ $ and
$ rk(G_{\alpha}(V))=d\alpha+s $.\medskip

Moreover, if $C$ has maximal rank, $C$ is minimally generated
if and only if \medskip
 
$rk(M_{\alpha+1}(V))= Min\{(r+1)($$\alpha+r\choose r$$-d\alpha-s) , 
$$\alpha+r+1\choose r$$-d(\alpha+1)-s)\}$\medskip \endproclaim
\demo{Proof} It follows immediately from Theorem 2.4 and (5),
(7)\qed\enddemo 
\proclaim {Corollary 2.6}  Let $C=$$\bigcup_{i=1}^s$$ C_i  \subset \Bbb
P_k^r$
 be a  non-degenerate  curve
such that $C_i$ are rational irreducible curves parametrized by maps 
$\Phi_i:\Bbb P_k^1 \rightarrow \Bbb P_k^r$,
 $\Phi_i([t_1,t_2])=[f_{i0}(t_1,t_2),...,f_{ir}(t_1,t_2)]$, 
where $f_{ij}\in k[t_1,t_2]$
are homogeneous polynomials of the same degree $d_i$. Let $a_1,...,a_q$
 $(q=s(d+1)(r+1))$
 be all the coefficients of all the polynomial $f_{ij}(t_1,t_2)$.
Then:\medskip
(a) the set $U\subset\Bbb A^q$ of  the $q$-tuples $(a_1,...,a_q)$ for
which 
 the curves 
$C_i$ are smooth of degree $d_i$ and disjoint, for any $i$, and $C$ has
maximal rank
is open;\medskip
(b) the  set $U'\subset\Bbb A^q$,
of the $q$-tuples $(a_1,...,a_q)\in U$ for which $C$ is minimally
generated is open.\endproclaim
\demo{Proof} By construction (see also Remark 2.3) the matrices 
$G_{\alpha-1}(V)$, $G_\alpha(V)$ and $M_{\alpha+1}(V)$ 
have entries which are forms in 
$a_1,...,a_q$ and then their  minors are forms in 
$a_1,...,a_q$. Moreover, by Proposition 1.3,(d) and (1),(5),(7), we have
$ rk(G_{\alpha-1}(V)) \leq c=$$\alpha+r-1\choose r$$ $,
$ rk(G_{\alpha}(V))\leq c'=d\alpha+s $ and 
$rk(M_{\alpha+1}(V)) \leq c"=Min\{(r+1)($$\alpha+r\choose r$$-d\alpha-s) , 
$$\alpha+r+1\choose r$$-d(\alpha+1)-s)\}$. Let  $M_1,...,M_u$  
 be  the  minors of 
$G_{\alpha-1}(V)$ of order $c$, $M'_1,...,M'_v$ be the minors of
$G_\alpha(V)$ of
order $c'$ and  $M"_1,...,M"_w$ be the minors of $M_{\alpha+1}(V)$ of
order  $c"$.
 Consider the 
following closed sets of $\Bbb A^q$, $D: M_1=0,...,M_u=0$, 
$D': M'_1=0,...,M'_v=0$,
$D": M"_1=0,...,M"_w=0$ and let $U=\Bbb A_k^q- (D\cap D')$ and
$U'=\Bbb A_k^q- (D\cap D'\cap D")$.  
If $(a_1,...,a_q)\in U$ then the matrices 
$G_{\alpha-1}(V)$ and $G_\alpha(V)$ have respectively rank $c$ and $c'$
and then $C$ has maximal rank 
[Corollary 2.5].
If $(a_1,...,a_q)\in U'$ then  the matrix $M_{\alpha+1}(V)$  has  rank
$c"$
and then $C$ has maximal rank and is minimally
generated [Corollary 2.5].\qed\enddemo
\proclaim{Remark 2.7} For some fixed
$r,s,d_1,...,d_s$ the open sets $U$ and $U'$ of Corollary $2.6$  could be
empty. For example  the union $C$ in $\Bbb P^3$ of
 two conics $C_1\cup C_2$ is contained
 in a reducible quadric and $C$ has not maximal rank. The union of $4$
generic 
lines in $\Bbb P^3$ has maximal rank but it is non minimally generated
$[9]$.\endproclaim 

\head 3.The ideal generation conjecture.\endhead
In this section we assume that the algebraically closed field $k$ has
characteristic zero.
 
Denote with $C=$$\bigcup_{i=1}^s$$ C_i  \subset \Bbb P^r$
 the union   of  rational irreducible smooth disjoint curves $C_i$.

Consider the following two questions:\medskip
 (a) {\sl For what values of $r, s, d_1,...,d_s$ is it true that, for
general
$C_1,...,C_s$, the curve $C=$$\bigcup_{i=1}^s$$ C_i  \subset \Bbb P^r$ has
maximal rank?}\medskip
(b) {\sl For what values of $r,n,d_1,...,d_s$ is it true that, for general
$C_1,...,C_s$, the curve $C=$$\bigcup_{i=1}^s$$ C_i  \subset \Bbb P^r$
 has maximal  rank and is finitely generated?}

Clearly, by Corollary 2.6, to answer to one of the two
questions for specific values of 
$r, s, d_1,...,d_s$,
it is enough to
find an example of curve $C$ with maximal rank or with maximal rank
 and finitely generated.
Question $(a)$  had a complete answer with the following result.
\proclaim{Theorem 3.1} Let $C=$$\bigcup_{i=1}^s$ $ C_i  \subset \Bbb
P_k^r$, $r\geq 3$,
be the union of $s$ rational irreducible smooth disjoint curves. Then, for
general
$C_1,...,C_s$, $C$ has maximal rank, for any $r,s,d_1,...,d_s$, 
 except the following $s$-uples $(d_1,...,d_s)$ of $\Bbb P_k^3$:\medskip
$(2,2);$ $(2,4);$ $(2,2,2);$ $(2,2,4);$ $(2,2,2,2)$.\endproclaim
\demo {Proof} If $r\geq 4$, for general $C_1,...,C_s$ any $C$ has maximal
rank  
[4, Theorem 1]. 

If $r=3$ the statement is proved in [1, Theorem 1]\qed\enddemo
Question $(b)$ has had a partial answer for lines and recently
for irreducible curves in $\Bbb P^3$ with the following results:
\proclaim{Theorem 3.2} (a) Let $C$
be the union of $s$   general lines in $\Bbb P^3$. Then $C$ has maximal
rank
 for any $s\neq 4$.

(b) If $C$ is a general irreducible rational curve of $\Bbb P^3$ of degree
$d\neq 5$
then
$C$ has maximal rank and is minimally generated if $d\leq 73$ or $d\geq
932$.\endproclaim
\demo{Proof} (a) is proved in [9, Theorem 1]

(b) has been recently proved in [3, Theorem 0.2 and Remark 8.139]  and in
[10, Theorem 1]
\qed\enddemo
In the following we want to show how the computer can be used to detect
 when a curve $C$ is of maximal rank and minimally generated.

Let $\Bbb Z_p$ be the residue field of the integers modulo a prime $p$.
If $\overline f=\overline b_1T_1+...+\overline b_uT_u$ is a polynomial
 of $\Bbb Z_p[t_1,t_2]$
($T_i$ are the terms in $t_1,t_2$   and $\overline b_i\in\Bbb Z_p$
are the coefficents  of $\overline f$) then
$f=b_1T_1+...+b_uT_u$ 
is the corresponding polynomial in $\Bbb Z[t_1,t_2]$. 
\proclaim{Theorem 3.3} 
Let $\overline C=\cup_i^s\overline C_i$ be the union of rational
irreducible curves $C_i\subset\Bbb P_{\Bbb Z_p^*}^r$ ($\Bbb Z_p^*$= 
algebraic closure of $\Bbb Z_p^*$)
 represented by forms $\overline f_{ij}\in \Bbb Z_p[t_1,t_2]$,
$j=1,...,r$.
Let $C=\cup_i^s C_i$ be the  union of the corresponding 
irreducible curves $C_i\subset\Bbb P_k^r$
 represented by the forms $f_{ij}\in \Bbb Z[t_1,t_2]\subset k[t_1,t_2]$,
  of degree $d_i$. Let $d=\sum_{i=1}^s d_i$ and 
$\alpha=Min\{n\in\Bbb N\mid $$n+r\choose r$$ > dn+s\}$
Suppose that  $\overline C_i$ contains $h_i=d_i(\alpha+1)+1$ distinct
points for any $i$.
let $\overline V$ be the set of all these $h=\sum_{i=1}^s
h_i=d(\alpha+1)+s$
 points. Consider the  matrices $(3)$ and $(6)$ of section $2$ relative to
the set $\overline V$.
If the following conditions hold:\medskip

(a) $ rk(G_{\alpha-1}(\overline V))=$$\alpha+r-1\choose r$$ $;\medskip
(b) 
$ rk(G_\alpha(\overline V))=d\alpha+s $;\medskip

(c) $rk(M_{\alpha+1}(\overline V))= Min\{(r+1)($$\alpha+r\choose
r$$-d\alpha-s),$$\alpha+r+1\choose r$$-d(\alpha+1)-s)\}$;\medskip

then 
  the curves $C_i$ are smooth and disjoint, $C$
has maximal rank and is minimally generated.\endproclaim
\demo{Proof} It is enough to recall the construction of the matrices
 given in section 2 and to observe the following: if $Q=(a_{ij})$ is a
matrix
  $m\times n$ with entries $a_{ij}\in \Bbb Z$
and $\overline Q= \overline a_{ij}$ is the matrix $m\times 
n$ whose entries are the integers
$a_{ij}$ modulo a prime $p$ then $rk (\overline Q)\leq rk (Q)$.
But, by Proposition 1.3,(d) and (1),(5),(7) of section 2, we have
 $ rk(G_{\alpha-1}( V)) \leq $$\alpha+r-1\choose r$$ $,
$ rk(G_\alpha( V)) \leq d\alpha+s $, and 

$rk(M_{\alpha+1}( V)) \leq  
Min\{(r+1)($$\alpha+r\choose r$$-d\alpha-s), 
$$\alpha+r+1\choose r$$-d(\alpha+1)-s)\}$.
  
 Hence, by the assumptions $(a)$,$(b)$,$(c)$, we have
 $ rk(G_{\alpha-1}(\overline V))=rk(G_{\alpha-1}( V))$,
$ rk(G_\alpha(\overline V))=rk(G_\alpha( V))$ and 
$ rk(M_{\alpha+1}(\overline V))=rk(M_{\alpha+1}( V))$ and the
claim follows from Corollary 2.5.\qed\enddemo


Then, fixed $r,s,d_1,...,d_s$, if we find an example of curve 
$\overline C=\cup_i^s \overline C_i$ 
for which the conditions $(a),(b),(c)$
of Theorem 3.3 are verified
 for a suitable prime $p$, then the corresponding curve
$C\subset\Bbb P_k^r$, 
 has maximal rank and
is minimally generated.

\proclaim{Remark 3.4} A key condition in Theorem $3.3$ is to find 
$h_i=d_i(\alpha+1)+1$ distinct points on any $\overline C_i$ but 
if $p$ is chosen large enough $($ e.g. $p=31991$ has worked in all our
examples$)$ 
$\overline C_i$ always contains $h_i$ distinct points $($ that are 
immediately found by giving
numerical values to the parameters $t_1,t_2$ that are used to represent
$C_i$,
 $i=1,...,s$ $)$.\endproclaim 

 Using Theorem 3.3 it is easy to costruct an algorithm that for fixed
$r,s,d_1,...,d_s$  produces examples (if they exist)
of curves $C$  with maximal rank and minimally generated 
\proclaim {Algorithm 3.5}\medskip
{\bf Step 1}. Input the integers $r,s,d_1,...,d_s$, the prime $p$ and the
 the coefficients (in $\Bbb Z_p$)  of the polynomials 
$\overline f_{ij}$ that represent $\overline C_i$.
\medskip
{\bf Step 2}. Find $h_i=d_i(\alpha+1)+1$ distinct points of $C_i$, for any
$i$
 (this can
be done, if $p$ is large enough, by giving numerical values to the
parameters
 $t_1,t_2$ of $\overline f_{ij}$).\medskip
{\bf Step 3}. Compute the ranks (with Gaussian reduction) of the matrices

$G_{\alpha-1}(\overline V)$,
 $G_\alpha(\overline V)$,
 $M_{\alpha+1}(\overline V)$.\medskip
{\bf Step 4}. If these three ranks are respectively
$ $$\alpha+r-1\choose r$$ $,
$  d\alpha+s $, and  
$Min\{(r+1)($$\alpha+r\choose r$$-d\alpha-s), 
$$\alpha+r+1\choose r$$-d(\alpha+1)-s)\}$ then output:
 $C$ has maximal rank and is minimally generated.
 
 
\endproclaim
\proclaim{Remark 3.6} By Corollary $2.6$ with a random choice of
coefficients in 
 $Step 1$ one always finds
a curve $C$ of maximal rank and finitely generated unless for fixed
$r,s,d_1,...,d_s$ there is $\underline {no}$ curve
 $C$ of maximal rank and finitely generated.\endproclaim
Using the language $C++$, G. Albano and F. Cioffi have implemented
Algorithm 3.5 
on our PC 
(an Intel Pentium $200$MHz with $40$ MB RAM$+$ $100$MB swap,
running Linux $2.0$).

The examples that we have run on the computer suggest the following: 
\proclaim{Conjecture 3.7 [Ideal Generation Conjecture ( I.G.C. ) for
rational curves]}
For every positive integers $r,s,d_1,...,d_s$, $r\geq 3$, 
the union $C=$$\bigcup_{i=1}^s$ $ C_i  \subset \Bbb P_k^r$ 
of  general rational irreducible  curves $C_i$ of degree $d_i$ has maximal
rank
and is minimally generated, 
 except for the following $s$-uples $(d_1,...,d_s)$.\medskip

If $r=3:$
$(5); (1,2); (2,2); (2,3); (2,4); (3,4); (1,1,2); (1,2,2); (2,2,2);
(2,2,3);$

$(2,2,4); (1,1,1,1); (1,2,2,2); (2,2,2,2);(1,2,2,2,2);$ \medskip

If $r=4:$
 $(5); (2,3); (1,1,2);$
\medskip
If $r=5:$
$(3,4); (2,2,2);$\medskip 
 \endproclaim
\proclaim{Remark 3.8} With geometric arguments it is easily seen that for
all the 
previous
exceptionals $s$-tuples all the corresponding curves have not maximal rank
or
have maximal rank but are not minimally generated.\endproclaim

In virtue of Corollary 2.6 to verify the conjecture it is enough to find
an example
for every $s$-uple $(d_1,...,d_s)$ different from the exceptional ones.
In particular we have verified by computer the following cases:
\proclaim{Theorem 3.9} Let $d=\sum_{i=1}^s d_i$ be the degree of $C$. If
$s=1$  then the I.G.C. is true for 
 $d\leq 80$. If $s\leq 10$  then the I.G.C. is true for 
 $d\leq 40$.\endproclaim

\proclaim{Remark 3.10} Theoretically there is no limit on the degree $d$
for  producing examples of union of rational curves
which satisfy the I.G.C. but when $s$ and $d$ get larger the computational
time can be very 
high and the cases to treat very numerous.\endproclaim
\proclaim {Remark 3.11} Curves $C=\cup_{i=1}^s  C_i$  which satisfy the
$I.G.C.$  could also be 
found by computing  a minimal set of generators of $I(C)$ 
with Computer Algebra Systems based on Groebner Bases computation, like
CoCoA. 
Using our PC, we  have computed,
 with the last version of CoCoA $[5]$ $($that uses for the elimination of
parameters
th Hilbert Driven Algorithm of Traverso$)$,
 a set of minimal generators of  $I(C)$ for general rational curves
$C_i$.
With CoCoA one needs first to compute a minimal set of generators of the
ideal
of every curve $C_i$ and then has to compute the generators of the
intersection of all the 
ideals $I(C_i)$. This  always needs much more computational time then with
our algorithm.
For example the minimal generators of the ideal of a general
 irreducible curve of degree $40$ in 
$\Bbb P^{15}$ are found with CoCoA in  $222$ 
 seconds and with our algorithm in $88$ seconds. 
If the irreducible curve has degree $80$ in $\Bbb P^{10}$ we have found
the 
minimal generators in  $310$ seconds while with CoCoA we did not get any
answer after
$9$ hours $($then we have turned off the computer$)$. 
The generators  of 
the ideal of $10$
curves with degree $8$ in $\Bbb P^{10}$ are computed with CoCoA in $3244$
seconds and
with our algorithm in $367$ seconds. Since to verify
 Theorem $3.8$
on the computer one needs to check hundreds of cases 
it is clear that CoCoA
 is less  suitable than our algorithm for checking I.G.C..
\endproclaim


 

  

 
\Refs
\widestnumber\key{14}
\ref\key 1
\by E. Ballico
\paper On the postulation of disjoint rational curves in a projective
space
\yr 1986 
\jour Rend. Sem. Mat. Univer. Politecn. Torino
\vol 44
\pages 46-52
\endref
\ref\key 2
\by E. Ballico
\paper Generators for the homogeneous ideal of $s$ general points in $\Bbb
P^3$
\yr 1987 
\jour J. Algebra
\vol 106
\pages 46-52
\endref

\ref\key 3
\by E. Ballico
\paper On the minimal free resolution of non-special curves in $\Bbb P^3$
\jour Preprint
\endref
\ref\key 4
\by E. Ballico, P. Ellia
\paper On the Postulation of Many Disjoint Rational Curves in $\Bbb P^n$,
$n\geq 4$                                                                  	
\yr 1985
\jour Bollettino U.M.I. 
\vol (6) 4-B
\pages 585-599
\endref
\ref\key 5
\by D. Capani, G. Niesi and L. Robbiano
\paper CoCoA (Computational Commutative Algebra)                                                                  	
\yr 1999
\jour available at http://lancelot.dima.unige.it 
\endref

\ref\key 6
\by A.V. Geramita, F. Orecchia
\paper Minimally Generating Ideal Defining Certain Tangent Cones                                                                   	
\yr 1982
\jour J. Algebra 
\vol 78
\pages 36-57
\endref
\ref\key 7
\by L. Gruson, R. Lazarsfeld, and C. Peskine
\paper On a Theorem of Castelnuovo, and the Equations Defining Space
Curves
\yr 1983
\jour  Invent. Math.
\vol 72
\pages 491-506
\endref
\ref\key 8
\by A. Hirschowitz, C. Simpson
\paper La r\'esolution minimale de l'id\'eal d'un arrangement g\'en\'eral
d'un grand nombre de
points dans $\Bbb P^n$                                                                   	
\yr 1996
\jour Invent. Math. 
\vol 126
\pages 467-503
\endref
\ref\key 9
\by M. Ida'
\paper On the homogeneous ideal of the generic union of $s$-lines in $\Bbb
P^3$ 
\yr 1990
\jour  J. Reine Angew. Math.
\vol 403
\pages 67-153
\endref
\ref\key 10
\by M. Ida'
\paper Generators for the generic rational space curve: low degree case 
\yr 1997
\jour  Proceedings of the Ferrara Conference in honour of M. Fiorentini
\endref

\ref\key 11
\by D. Mumford
\paper Lectures on curves on an algebraic surface
\yr 1966
\jour Ann. of Math. studies
\vol 59
\endref
\endRefs
\enddocument